\chapter{2009年福建卷（理）}

\section{单选题}

\begin{problem}
函数 \( f(x) = \sin x\cos x \) 的最小值是（\quad）

\choices
    {\(-1\)}
    {\(-\frac{1}{2}\)}
    {\( \frac{1}{2} \)}
    {\(1\)}
\answer{B}

\begin{solution}
由二倍角公式，\(f(x)=\sin x\cos x=\frac{1}{2}\sin 2x\).因为\(-1\leqslant\sin 2x\leqslant1\)，所以
\(f(x)\geqslant-\frac{1}{2}\)，且当\(\sin 2x=-1\)时取等号.因此最小值为\(-\frac{1}{2}\)，选 B.
\end{solution}
\end{problem}

\begin{problem}
已知全集 \(U = \R\)，集合 \(A = \{x \mid x^2 - 2x > 0\}\)，则 \(\complement_U A\) 等于（\quad）

\choices
    {\(\{x\mid 0\leqslant x\leqslant 2\}\)}
    {\(\{x\mid 0 < x < 2\}\)}
    {\(\{x\mid x < 0\text{ 或 }x > 2\}\)}
    {\(\{x\mid x\leqslant 0\text{ 或 }x\geqslant 2\}\)}
\answer{A}

\begin{solution}
由\(x^2-2x=x(x-2)>0\)，得\(x<0\)或\(x>2\)，所以
\(A=(-\infty,0)\cup(2,+\infty)\).在全集\(\R\)中取补集，得到
\(\complement_U A=[0,2]=\{x\mid 0\leqslant x\leqslant2\}\)，选 A.
\end{solution}
\end{problem}

\begin{problem}
等差数列 \(\{a_{n}\}\) 的前 \(n\) 项和为 \(S_{n}\)，且 \(S_{3} = 6\)，\(a_{1} = 4\)，则公差 \(d\) 等于（\quad）

\choices
    {\(1\)}
    {\( \frac{5}{3} \)}
    {\(2\)}
    {\(3\)}
\answer{C}

\begin{solution}
等差数列前三项为\(a_1=4\)，\(a_2=4+d\)，\(a_3=4+2d\).由题设
\[
S_3=a_1+a_2+a_3=4+(4+d)+(4+2d)=12+3d=6
\]
解得\(d=-2\).因此按题设条件公差为\(-2\)；原 PDF 的选项 C 印为\(2\)，与计算结果不一致.
\end{solution}
\end{problem}

\begin{problem}
\(\int_{-\frac{\pi}{2}}^{\frac{\pi}{2}}(1 + \cos x)\mathrm{d}x\) 等于（\quad）

\choices
    {\(\pi\)}
    {\(2\)}
    {\(\pi - 2\)}
    {\(\pi + 2\)}
\answer{D}

\begin{solution}
由原函数的一个原函数可得
\[
\int_{-\frac{\pi}{2}}^{\frac{\pi}{2}}(1+\cos x)\mathrm{d}x
=\left.x+\sin x\right|_{-\frac{\pi}{2}}^{\frac{\pi}{2}}
=\left(\frac{\pi}{2}+1\right)-\left(-\frac{\pi}{2}-1\right)=\pi+2
\]
因此选 D.
\end{solution}
\end{problem}

\begin{problem}
下列函数 \( f(x) \) 中，满足“对任意 \(x_{1}\)，\(x_{2} \in (0, +\infty)\)，当 \( x_{1} < x_{2} \) 时，都有 \( f(x_{1}) > f(x_{2}) \)”的是（\quad）

\choices
    {\(f(x) = \frac{1}{x}\)}
    {\(f(x) = (x - 1)^{2}\)}
    {\(f(x) = \e^{x}\)}
    {\(f(x) = \ln (x + 1)\)}
\answer{A}

\begin{solution}
当\(x>0\)时，\(f(x)=\frac{1}{x}\)的导数为\(f'(x)=-\frac{1}{x^2}<0\)，所以它在\((0,+\infty)\)上严格递减，满足题设.
\(f(x)=(x-1)^2\)在\((0,1)\)上递减、在\((1,+\infty)\)上递增；\(f(x)=\e^x\)和\(f(x)=\ln(x+1)\)的导数分别为\(\e^x>0\)和\(\frac{1}{x+1}>0\)，均严格递增.因此选 A.
\end{solution}
\end{problem}

\begin{problem}
阅读如图所示的程序框图，运行相应的程序，输出的结果是（\quad）

\bitmapfigure[max width=0.42\linewidth]{img/2009/fujian_science/q6_fig1.png}

\choices
    {\(2\)}
    {\(4\)}
    {\(8\)}
    {\(16\)}
\answer{C}

\begin{solution}
程序开始时\(S=2\)，\(n=1\).循环第一次计算
\(S=\frac{1}{1-S}=\frac{1}{1-2}=-1\)，随后\(n=2\)；第二次计算
\(S=\frac{1}{1-(-1)}=\frac{1}{2}\)，随后\(n=4\)；第三次计算
\(S=\frac{1}{1-\frac{1}{2}}=2\)，随后\(n=8\).此时满足\(S=2\)，输出\(n=8\)，选 C.
\end{solution}
\end{problem}

\begin{problem}
设 \(m\)，\(n\) 是平面 \(\alpha\) 内的两条不同直线，\(l_{1}\)，\(l_{2}\) 是平面 \(\beta\) 内的两条相交直线，则 \(\alpha\parallel\beta\) 的一个充分而不必要条件是（\quad）

\choices
    {\(m\parallel\beta\) 且 \(l_{1}\parallel\alpha\)}
    {\(m\parallel l_{1}\) 且 \(n\parallel l_{2}\)}
    {\(m\parallel\beta\) 且 \(n\parallel\beta\)}
    {\(m\parallel\beta\) 且 \(n\parallel l_{2}\)}
\answer{B}

\begin{solution}
若\(m\parallel l_1\)且\(n\parallel l_2\)，则\(m\)与\(n\)的方向分别平行于平面\(\beta\)内相交直线\(l_1\)与\(l_2\)的方向.由于\(m\)，\(n\)在平面\(\alpha\)内且方向不同，平面\(\alpha\)的两个方向均与平面\(\beta\)的两个方向平行，从而\(\alpha\)与\(\beta\)平行.因此 B 是充分条件.

它不是必要条件.例如两平面平行时，可以在\(\alpha\)内任取两条不同直线\(m\)，\(n\)，它们未必分别平行于指定的\(l_1\)，\(l_2\).所以该条件充分而不必要，选 B.
\end{solution}
\end{problem}

\begin{problem}
已知某运动员每次投篮命中的概率都为 \(40\%\). 现采用随机模拟的方法估计该运动员三次投篮恰有两次命中的概率：先由计算器算出\(0\)到\(9\)之间取整数值的随机数，指定\(1\)，\(2\)，\(3\)，\(4\)表示命中，\(5\)，\(6\)，\(7\)，\(8\)，\(9\)，\(0\)表示不命中；再以每\(3\)个随机数为一组，代表三次投篮的结果. 经随机模拟产生了\(20\)组随机数：

\begin{center}
\begin{tabular}{ccccc}
\(907\) & \(966\) & \(191\) & \(925\) & \(271\) \\
\(932\) & \(812\) & \(458\) & \(569\) & \(683\) \\
\(431\) & \(257\) & \(393\) & \(027\) & \(556\) \\
\(488\) & \(730\) & \(113\) & \(537\) & \(989\)
\end{tabular}
\end{center}

据此估计，该运动员三次投篮有两次命中的概率为（\quad）

\choices
    {\(0.35\)}
    {\(0.25\)}
    {\(0.20\)}
    {\(0.15\)}
\answer{B}

\begin{solution}
把每组的三个数字按“\(1\)，\(2\)，\(3\)，\(4\) 命中，其余不命中”判断.恰有两次命中的组为
\(191\)、\(271\)、\(932\)、\(812\)、\(393\)，共\(5\)组.总共模拟了\(20\)组，所以频率为
\(\frac{5}{20}=0.25\).因此估计的概率为\(0.25\)，选 B.
\end{solution}
\end{problem}

\begin{problem}
设 \(\bs{a}\)，\(\bs{b}\)，\(\bs{c}\) 为同一平面内具有相同起点的任意三个非零向量，且满足 \(\bs{a}\) 与 \(\bs{b}\) 不共线，\(\bs{a} \perp \bs{c}\)，\(|\bs{a}| = |\bs{c}|\)，则 \(\lvert \bs{b} \cdot \bs{c} \rvert\) 的值一定等于（\quad）

\choices
    {以 \(\bs{a}, \bs{b}\) 为两边的三角形面积}
    {以 \(\bs{b}\)，\(\bs{c}\) 为两边的三角形面积}
    {以 \(\bs{a}\)，\(\bs{b}\) 为邻边的平行四边形的面积}
    {以 \(\bs{b}\)，\(\bs{c}\) 为邻边的平行四边形的面积}
\answer{C}

\begin{solution}
设\(\theta\)为\(\bs{a}\)与\(\bs{b}\)的夹角.因为\(\bs{a}\perp\bs{c}\)且\(|\bs{a}|=|\bs{c}|\)，所以\(\bs{b}\)与\(\bs{c}\)的夹角的余弦绝对值等于\(\sin\theta\).于是
\[
\lvert\bs{b}\cdot\bs{c}\rvert
=|\bs{b}||\bs{c}|\lvert\cos\angle(\bs{b},\bs{c})\rvert
=|\bs{b}||\bs{a}|\sin\theta
\]
右端正是以\(\bs a\)，\(\bs b\)为邻边的平行四边形面积.因此选 C.
\end{solution}
\end{problem}

\begin{problem}
函数 \( f(x) = ax^2 + bx + c \)（\(a \neq 0\)）的图象关于直线 \( x = -\frac{b}{2a} \) 对称. 据此可推测，对任意的非零实数 \(a\)，\(b\)，\(c\)，\(m\)，\(n\)，\(p\)，关于 \( x \) 的方程 \( m[f(x)]^2 + nf(x) + p = 0 \) 的解集都不可能是（\quad）

\choices
    {\( \{1, 2\} \)}
    {\( \{1, 4\} \)}
    {\( \{1, 2, 3, 4\} \)}
    {\( \{1, 4, 16, 64\} \)}
\answer{D}

\begin{solution}
二次函数图象关于直线\(x=-\frac{b}{2a}\)对称，所以对任意使\(f(x)\)有定义的\(x\)，关于该直线对称的点
\(x'=-\frac{b}{a}-x\)满足\(f(x')=f(x)\).因此方程
\(m[f(x)]^2+nf(x)+p=0\)的解集必须关于同一条竖直直线对称：方程先确定至多两个\(f(x)\)的取值，而每个取值的原像关于抛物线的对称轴成对出现.

集合\(\{1,2\}\)、\(\{1,4\}\)分别关于\(x=\frac{3}{2}\)、\(x=\frac{5}{2}\)对称，集合\(\{1,2,3,4\}\)关于\(x=\frac{5}{2}\)对称，均不违反这个必要条件.它们确实可以实现. 对于集合\(\{1,2\}\)，取\(f(x)=\left(x-\frac{3}{2}\right)^2+1\)，并令关于\(y\)的二次方程为\((y-\tfrac{5}{4})(y+1)=0\)；对于集合\(\{1,4\}\)，取\(f(x)=\left(x-\frac{5}{2}\right)^2+1\)，并令关于\(y\)的二次方程为\((y-\tfrac{13}{4})(y+1)=0\)；对于集合\(\{1,2,3,4\}\)，取\(f(x)=\left(x-\frac{5}{2}\right)^2\)，并令关于\(y\)的二次方程为\((y-\tfrac{9}{4})(y-\tfrac{1}{4})=0\).
前三个例子展开后\(a\)，\(b\)，\(c\)，\(m\)，\(n\)，\(p\)均非零.第一个方程的解集为\(\{1,2\}\)，第二个方程的解集为\(\{1,4\}\)，第三个方程的解集为\(\{1,2,3,4\}\).而\(\{1,4,16,64\}\)中，若关于同一条竖直直线成对对称，则应有\(1+64=4+16\)，但\(65\neq20\)，故该集合不可能是解集.因此不可能的是 D.
\end{solution}
\end{problem}

\section{填空题}

\begin{problem}
若 \( \frac{2}{1 - \i} = a + b\i \)（\(\i\) 为虚数单位，\(a\)，\(b\in\R\)），则 \(a+b=\)\fillinblank{}.
\answer{\(2\)}

\begin{solution}
分母有理化得
\[
\frac{2}{1-\i}=\frac{2(1+\i)}{(1-\i)(1+\i)}=\frac{2(1+\i)}{2}=1+\i
\]
所以\(a=1\)，\(b=1\)，从而\(a+b=2\).
\end{solution}
\end{problem}

\begin{problem}
某校开展“爱我海西，爱我家乡”摄影比赛，\(9\)位评委为参赛作品 \(A\) 给出的分数如茎叶图所示. 记分员在去掉一个最高分和一个最低分后，算术平均分为\(91\)，复核员在复核时，发现有一个数字（茎叶图中的 \(x\)）无法看清. 若记分员计算无误，则数字 \(x\) 应该是\fillinblank{}.

\begin{center}
\begin{tabular}{c|cccccc}
\phantom{8} & \multicolumn{6}{|c}{作品 A} \\
\hline
\(8\) & \(8\) & \(9\) & \(9\) & & & \\
\(9\) & \(2\) & \(3\) & \(x\) & \(2\) & \(1\) & \(4\)
\end{tabular}
\end{center}
\answer{\(1\)}

\begin{solution}
茎叶图中的九个分数为\(88\)，\(89\)，\(89\)，\(91\)，\(92\)，\(92\)，\(93\)，\(94\)，\(90+x\).去掉最低分\(88\)后，最高分为\(94\)（当\(x\leqslant4\)）或\(90+x\)（当\(x\geqslant5\)）.去掉最高、最低分后，七个分数的和应为\(7\times91=637\).

九个分数的总和为\(818+x\).若\(x\leqslant4\)，去掉\(88,94\)后有
\(818+x-88-94=637\)，解得\(x=1\)，符合\(x\leqslant4\).若\(x\geqslant5\)，去掉\(88\)，\(90+x\)后总和为\(640\)，不等于\(637\).因此\(x=1\).
\end{solution}
\end{problem}

\begin{problem}
过抛物线 \( y^{2} = 2px  \) （\(p > 0\)）的焦点 \(F\) 作倾斜角为 \( 45^{\circ} \) 的直线交抛物线于 \(A\)，\(B\) 两点，若线段 \(AB\) 的长为 \( 8 \)，则 \( p = \)\fillinblank{}.
\answer{\(2\)}

\begin{solution}
抛物线\(y^2=2px\)的焦点为\(F\left(\frac{p}{2},0\right)\).倾斜角为\(45^{\circ}\)的直线为
\(y=x-\frac{p}{2}\).代入抛物线得
\[
\left(x-\frac{p}{2}\right)^2=2px
\quad\Longleftrightarrow\quad
x^2-3px+\frac{p^2}{4}=0
\]
设两交点的横坐标为\(x_1\)，\(x_2\)，则
\(\lvert x_1-x_2\rvert=\sqrt{9p^2-p^2}=2\sqrt{2}p\).由于直线斜率为\(1\)，弦长
\[
AB=\sqrt{(x_1-x_2)^2+(y_1-y_2)^2}
=\sqrt{2}\lvert x_1-x_2\rvert=4p
\]
由\(AB=8\)得\(p=2\).
\end{solution}
\end{problem}

\begin{problem}
若曲线 \( f(x) = ax^3 + \ln x \) 存在垂直于 \( y \) 轴的切线，则实数 \( a \) 取值范围是\fillinblank{}.
\answer{\((-\infty,0)\)}

\begin{solution}
函数定义域为\(x>0\)，且
\(f'(x)=3ax^2+\frac{1}{x}\).切线垂直于\(y\)轴即切线水平，所以存在\(x_0>0\)使
\[
f'(x_0)=0
\quad\Longleftrightarrow\quad
3ax_0^3+1=0
\quad\Longleftrightarrow\quad
a=-\frac{1}{3x_0^3}<0
\]
反之，任取\(a<0\)，取\(x_0=\sqrt[3]{-\frac{1}{3a}}>0\)，即可使\(f'(x_0)=0\).因此\(a\)的取值范围是\((-\infty,0)\).
\end{solution}
\end{problem}

\begin{problem}
五位同学围成一圈依序循环报数，规定：

\begin{circlelist}
\item 第一位同学首次报出的数为 \(1\)，第二位同学首次报出的数也为 \(1\)，之后每位同学所报出的数都是前两位同学所报出的数之和；
\item 若报出的数为 \(3\) 的倍数，则报该数的同学需拍手一次.
\end{circlelist}

已知甲同学第一个报数，当五位同学依序循环报到第\(100\)个数时，甲同学拍手的总次数为\fillinblank{}.
\answer{\(5\)}

\begin{solution}
记第\(n\)个报出的数为\(F_n\)，则\(F_1=F_2=1\)，且\(F_n=F_{n-1}+F_{n-2}\).模\(3\)计算得到
\(F_1\)，\(F_2\)，\(\ldots\)，\(F_8\equiv1\)，\(1\)，\(2\)，\(0\)，\(2\)，\(2\)，\(1\)，\(0\pmod{3}\)
又 \(F_9\equiv F_1\)，\(F_{10}\equiv F_2\pmod{3}\)，递推关系相同，所以之后以\(8\)为周期.因此\(F_n\)是\(3\)的倍数当且仅当\(n\equiv0\)或\(4\pmod{8}\).

甲同学第一次报第\(1\)个数，五人循环，所以甲报数的序号为\(n=1+5k\).报到第\(100\)个数时\(0\leqslant k\leqslant19\).条件
\(1+5k\equiv0\)或\(4\pmod{8}\)分别给出\(k\equiv3\)或\(7\pmod{8}\).在\(0\leqslant k\leqslant19\)中，这些\(k\)为\(3\)，\(7\)，\(11\)，\(15\)，\(19\)，共\(5\)个.因此甲同学拍手\(5\)次.
\end{solution}
\end{problem}

\section{解答题}

\begin{problem}
从集合 \(\{1, 2, 3, 4, 5\}\) 的所有非空子集中，等可能地取出一个.

\begin{enumerate}
\item 记性质 \(r\)：集合中的所有元素之和为 \(10\)，求所取出的非空子集满足性质 \(r\) 的概率；
\item 记所取出的非空子集的元素个数为 \(\xi\)，求 \(\xi\) 的分布列和数学期望 \(E_{\xi}\).
\end{enumerate}

\begin{solution}
\begin{enumerate}
\item 非空子集共有\(2^5-1=31\)个，且每个被取出的可能性相同.元素和为\(10\)的非空子集为
\(\{1,4,5\}\)、\(\{2,3,5\}\)、\(\{1,2,3,4\}\)，共\(3\)个，所以\(P(r)=\dfrac{3}{31}\).

\item 记所取子集的元素个数为\(\xi\).当\(\xi=k\)时，从\(5\)个元素中选出\(k\)个，共有\(\mathrm C_5^k\)种，故
\begin{center}
\begin{tabular}{c|ccccc}
\(\xi\)&\(1\)&\(2\)&\(3\)&\(4\)&\(5\)\\
\hline
\(P\)&\(\dfrac{5}{31}\)&\(\dfrac{10}{31}\)&\(\dfrac{10}{31}\)&\(\dfrac{5}{31}\)&\(\dfrac{1}{31}\)
\end{tabular}
\end{center}
并且
\[
E_{\xi}=1\cdot\frac{5}{31}+2\cdot\frac{10}{31}+3\cdot\frac{10}{31}+4\cdot\frac{5}{31}+5\cdot\frac{1}{31}
=\frac{80}{31}
\]
\end{enumerate}
\end{solution}
\end{problem}

\begin{problem}
如图，四边形 \(ABCD\) 是边长为 1 的正方形，\(MD \perp\) 平面 \(ABCD\)，\(NB \perp\) 平面 \(ABCD\)，且 \(MD = NB = 1\)，\(E\) 为 \(BC\) 的中点.

\begin{enumerate}
\item 求异面直线 \(NE\) 与 \(AM\) 所成角的余弦值；
\item 在线段 \(AN\) 上是否存在点 \(S\)，使得 \(ES \perp\) 平面 \(AMN\)？若存在，求线段 \(AS\) 的长；若不存在，请说明理由.
\end{enumerate}
\bitmapfigure[max width=0.42\linewidth]{img/2009/fujian_science/q17_fig1.png}

\begin{solution}
\begin{enumerate}
\item 建立空间直角坐标系，使\(D\)为原点，\(DC\)、\(DA\)、\(DM\)分别为三个坐标轴正方向.则\(A(0,1,0)\)、\(B(1,1,0)\)、\(C(1,0,0)\)、\(M(0,0,1)\)、\(N(1,1,1)\)、\(E\left(1,\frac{1}{2},0\right)\).
于是\(\overrightarrow{NE}=\left(0,-\frac{1}{2},-1\right)\)，\(\overrightarrow{AM}=(0,-1,1)\)，且
\(\overrightarrow{NE}\cdot\overrightarrow{AM}=-\frac{1}{2}\)、\(|\overrightarrow{NE}|=\frac{\sqrt{5}}{2}\)、\(|\overrightarrow{AM}|=\sqrt{2}\).
异面直线所成角取锐角，故其余弦值为
\[
\frac{|\overrightarrow{NE}\cdot\overrightarrow{AM}|}{|\overrightarrow{NE}||\overrightarrow{AM}|}
=\frac{\frac{1}{2}}{\frac{\sqrt{10}}{2}}=\frac{\sqrt{10}}{10}
\]

\item 设\(S\)在线段\(AN\)上，令\(S=A+t(N-A)=(t,1,t)\)，其中\(0\leqslant t\leqslant1\).平面\(AMN\)内有
\(\overrightarrow{AM}=(0,-1,1)\)，\(\overrightarrow{AN}=(1,0,1)\)，其法向量可取
\[
\overrightarrow{AM}\times\overrightarrow{AN}=(-1,1,1)
\]
又\(\overrightarrow{ES}=\left(t-1,\frac{1}{2},t\right)\).要使\(ES\perp\)平面\(AMN\)，必须有\(\overrightarrow{ES}\)平行于\((-1,1,1)\).比较第二坐标得比例因子为\(\frac{1}{2}\)，于是第一、第三坐标分别给出\(t-1=-\frac{1}{2}\)和\(t=\frac{1}{2}\)，两者一致.因此存在这样的点\(S\)，且
\[
AS=t|AN|=\frac{1}{2}\sqrt{1^2+1^2}=\frac{\sqrt{2}}{2}
\]
\end{enumerate}
\end{solution}
\end{problem}

\begin{problem}
如图，某市拟在长为 \(8 \mathrm{~km}\) 的道路 \(OP\) 的一侧修建一条运动赛道. 赛道的前一部分为曲线段 \(OSM\)，该曲线段为函数 \(y = A \sin \omega x\)（\(A > 0\)，\(\omega > 0\)），\(x \in [0,4]\) 的图象，且图象的最高点为 \(S(3, 2\sqrt{3})\)；赛道的后一部分为折线段 \(MNP\). 为保证参赛运动员的安全，限定 \(\angle MNP = 120^{\circ}\).

\begin{enumerate}
\item 求 \(A\)，\(\omega\) 的值和 \(M\)，\(P\) 两点间的距离；
\item 应如何设计，才能使折线段赛道 \(MNP\) 最长？
\end{enumerate}
\bitmapfigure[max width=0.48\linewidth]{img/2009/fujian_science/q18_fig1.png}

\begin{solution}
\begin{enumerate}
\item 因为图象的最高点为\(S(3,2\sqrt{3})\)，所以振幅\(A=2\sqrt{3}\)，且\(\sin(3\omega)=1\).写成\(3\omega=\frac{\pi}{2}+2k\pi\).由\(\omega>0\)知\(k\geqslant0\).若\(k\geqslant1\)，则前一个最高点的横坐标为\(3-\frac{2\pi}{\omega}\in(0,3)\)，与最高点只有\(S\)不符，故\(k=0\)，从而\(\omega=\frac{\pi}{6}\).

曲线在\(x=4\)处的终点为
\[
M\left(4,2\sqrt{3}\sin\frac{2\pi}{3}\right)=(4,3)
\]
而\(P=(8,0)\)，所以
\[
MP=\sqrt{(8-4)^2+(0-3)^2}=5
\]

\item 设\(MN=u\)，\(NP=v\)，其中\(u\)，\(v>0\).在三角形\(MNP\)中，\(\angle MNP=120^{\circ}\)，由余弦定理
\[
MP^2=u^2+v^2-2uv\cos120^{\circ}=u^2+v^2+uv=25
\]
因此由\(u^2+v^2\geqslant2uv\)得\(25\geqslant3uv\)，即\(uv\leqslant\frac{25}{3}\).于是
\[
(u+v)^2=u^2+v^2+2uv=25+uv\leqslant25+\frac{25}{3}=\frac{100}{3}
\]
等号当且仅当\(u=v\)，即\(MN=NP\).此时三角形为等腰三角形，底角
\(\angle PMN=\angle MPN=\frac{180^{\circ}-120^{\circ}}{2}=30^{\circ}\).所以应设计为\(MN=NP\)，等价于\(\angle PMN=30^{\circ}\).
\end{enumerate}
\end{solution}
\end{problem}

\begin{problem}
已知 \(A\)，\(B\) 分别为曲线 \(C: \frac{x^2}{a^2} + y^2 = 1\)（\(y \geqslant 0\)，\(a > 0\)）与 \(x\) 轴的左，右两个交点，直线 \(l\) 过点 \(B\)，且与 \(x\) 轴垂直，\(S\) 为 \(l\) 上异于点 \(B\) 的一点，连结 \(AS\) 交曲线 \(C\) 于点 \(T\).

\begin{enumerate}
\item 若曲线 \(C\) 为半圆，点 \(T\) 为圆弧 \( \widehat{AB} \) 的三等分点，试求出点 \(S\) 的坐标；
\item 如图，点 \(M\) 是以 \(SB\) 为直径的圆与线段 \(TB\) 的交点. 试问：是否存在 \(a\)，使得 \(O\)，\(M\)，\(S\) 三点共线？若存在，求出 \(a\) 的值，若不存在，请说明理由.
\end{enumerate}
\bitmapfigure[max width=0.48\linewidth]{img/2009/fujian_science/q19_fig1.png}

\begin{solution}
\begin{enumerate}
\item 当\(C\)为半圆时\(a=1\)，故\(A=(-1,0)\)，\(B=(1,0)\).圆弧\(AB\)的两个三等分点为\(T_1\left(\frac{1}{2},\frac{\sqrt{3}}{2}\right)\)和\(T_2\left(-\frac{1}{2},\frac{\sqrt{3}}{2}\right)\).

直线\(AT_1\)的斜率为\(\frac{\sqrt{3}}{3}\)，代入\(x=1\)得\(y=\frac{2\sqrt{3}}{3}\)；直线\(AT_2\)的斜率为\(\sqrt{3}\)，代入\(x=1\)得\(y=2\sqrt{3}\).因此
\[
S=\left(1,\frac{2\sqrt{3}}{3}\right)\quad\text{或}\quad S=(1,2\sqrt{3})
\]

\item 设\(S=(a,s)\).由于第二个交点\(T\)在\(y\geqslant0\)的上半椭圆上，且\(S\neq B\)，故取\(s>0\).令直线\(AS\)上的点为\((x,y)=(-a+2au,su)\).代入椭圆方程，除去对应于\(A\)的根\(u=0\)，得到
\((-1+2u)^2+s^2u^2=1\)，\(\Longrightarrow\)，\(u=\frac{4}{4+s^2}\)
所以
\[
T=\left(a\left(-1+\frac{8}{4+s^2}\right),\frac{4s}{4+s^2}\right)
\]
圆\(SB\)以\(SB\)为直径，其方程为\((x-a)^2+y^2-sy=0\).若\(O\)，\(M\)，\(S\)共线，设\(M=\lambda S=(\lambda a,\lambda s)\).由于\(M\)在\(TB\)上而\(S\)不在\(TB\)上，故\(\lambda\neq1\).由\(\angle SMB=90^{\circ}\)得
\[
(S-M)\cdot(B-M)=(1-\lambda)\bigl[a^2(1-\lambda)-\lambda s^2\bigr]=0
\]
由于\(\lambda\neq1\)，所以
\[
a^2(1-\lambda)-\lambda s^2=0
\]
从而\(\lambda=\frac{a^2}{a^2+s^2}\)，即
\[
M=\left(\frac{a^3}{a^2+s^2},\frac{a^2s}{a^2+s^2}\right)
\]
另一方面，由上式中\(T\)的坐标，直线\(BT\)的斜率为\(\frac{y_T}{x_T-a}=-\frac{2}{as}\)；而\(BM\)的斜率为\(-\frac{a}{s}\).要使\(B\)，\(T\)，\(M\)共线，必须有\(-\frac{a}{s}=-\frac{2}{as}\)，所以\(a^2=2\).因为\(a>0\)，得到\(a=\sqrt{2}\).此时
\(M-B=\frac{4+s^2}{2(2+s^2)}(T-B)\)，\(0<\frac{4+s^2}{2(2+s^2)}<1\)
故\(M\)确在线段\(TB\)上，所求值存在.
\end{enumerate}
\end{solution}
\end{problem}

\begin{problem}
已知函数 \( f(x) = \frac{1}{3} x^3 + ax^2 + bx \)，且 \( f'(-1) = 0 \).

\begin{enumerate}
\item 试用含 \(a\) 的代数式表示 \(b\)，并求 \(f(x)\) 的单调区间；
\item 令 \(a = -1\)，设函数 \(f(x)\) 在 \(x_1\)，\(x_2\)（\(x_1 < x_2\)）处取得极值，记点 \(M(x_1, f(x_1))\)，\(N(x_2, f(x_2))\)，\(P(m, f(m))\)，\(x_1 < m < x_2\). 请仔细观察曲线 \(f(x)\) 在点 \(P\) 处的切线与线段 \(MP\) 的位置变化趋势，并解释以下问题：

\begin{enumerate}
\item 若对任意的 \(m \in (t, x_2]\)，线段 \(MP\) 与曲线 \(f(x)\) 均有异于 \(M\)，\(P\) 的公共点，试确定 \(t\) 的最小值，并证明你的结论；
\item 若存在点 \(Q(n, f(n))\)，\(x_1 \leqslant n < m\)，使得线段 \(PQ\) 与曲线 \(f(x)\) 有异于 \(P\)，\(Q\) 的公共点，请直接写出 \(m\) 的取值范围. （不必给出求解过程）
\end{enumerate}
\end{enumerate}

\begin{solution}
\begin{enumerate}
\item 由
\(f'(x)=x^2+2ax+b\)，\(f'(-1)=1-2a+b=0\)
得\(b=2a-1\).于是
\[
f'(x)=x^2+2ax+2a-1=(x+1)(x+2a-1)=(x+1)(x-(1-2a))
\]

当\(a>1\)时，\(1-2a<-1\)，故\(f'(x)\)在\((-\infty,1-2a)\)和\((-1,+\infty)\)上为正，在\((1-2a,-1)\)上为负.因此增区间为\((-\infty,1-2a)\)和\((-1,+\infty)\)，减区间为\((1-2a,-1)\).

当\(a=1\)时，\(f'(x)=(x+1)^2\geqslant0\)，所以\(f(x)\)在\(\R\)上单调递增.

当\(a<1\)时，\(-1<1-2a\)，故\(f'(x)\)在\((-\infty,-1)\)和\((1-2a,+\infty)\)上为正，在\((-1,1-2a)\)上为负.因此增区间为\((-\infty,-1)\)和\((1-2a,+\infty)\)，减区间为\((-1,1-2a)\).

\item 令\(a=-1\)，则\(b=-3\)，故\(f(x)=\frac{1}{3}x^3-x^2-3x\)，\(f'(x)=(x+1)(x-3)\)，所以\(x_1=-1\)，\(x_2=3\).

对曲线上横坐标为\(u\)，\(v\)的两个点作割线.由于\(f(x)\)的三次项系数为\(\frac{1}{3}\)，割线方程与曲线方程相减后所得三次方程的三个根为\(u\)，\(v\)，\(w\).该三次方程的首项为\(\frac{1}{3}x^3\)，而二次项系数为\(-1\)，与\(\frac{1}{3}(x-u)(x-v)(x-w)\)比较可得\(-\frac{u+v+w}{3}=-1\)，即\(u+v+w=3\)，所以第三个交点的横坐标为\(w=3-u-v\).

\begin{enumerate}
\item 取\(u=x_1=-1\)，\(v=m\)，得\(w=4-m\).当\(m\in(t,3]\)时，线段\(MP\)与曲线有异于\(M\)，\(P\)的公共点，要求第三个交点的横坐标满足\(-1<w<m\).代入\(w=4-m\)得\(4-m<m\)，即\(m>2\)；当\(m=3\)时仍有\(w=1\)，满足条件.因此对所有\(m\in(2,3]\)均成立，而任意\(t<2\)都会使\(m=2\)落入\((t,3]\)，此时\(w=m\)并无异于\(P\)的第三个交点.故\(t\)的最小值为\(2\).

\item 取\(u=n\)，\(v=m\)，第三个交点横坐标为\(w=3-m-n\).它位于线段\(PQ\)内部当且仅当
\[
n<3-m-n<m
\quad\Longleftrightarrow\quad
3-2m<n<\frac{3-m}{2}
\]
还需满足\(-1\leqslant n<m\).当\(m\leqslant1\)时，\(3-2m\geqslant m\)，不存在这样的\(n\)；当\(1<m<3\)时，上述不等式区间与\([-1,m)\)有公共部分，可以选取这样的\(n\).所以\(m\)的取值范围为\((1,3)\).
\end{enumerate}
\end{enumerate}
\end{solution}
\end{problem}

\section{选考题：解答题}

\begin{problem}
已知矩阵 \( M = \begin{pmatrix} 2 & -3 \\ 1 & -1 \end{pmatrix} \) 所对应的线性变换把点 \( A(x, y) \) 变成点 \( A'(13, 5) \)，试求 \(M\) 的逆矩阵及点 \(A\) 的坐标.

\begin{solution}
矩阵\(M\)的行列式为
\(\det M=2\times(-1)-(-3)\times1=1\neq0\)，所以\(M\)可逆，且
\[
M^{-1}=\frac{1}{\det M}\begin{pmatrix}-1&3\\-1&2\end{pmatrix}
=\begin{pmatrix}-1&3\\-1&2\end{pmatrix}
\]
由线性变换关系
\[
M\begin{pmatrix}x\\y\end{pmatrix}=\begin{pmatrix}13\\5\end{pmatrix}
\]
得
\[
\begin{pmatrix}x\\y\end{pmatrix}=M^{-1}\begin{pmatrix}13\\5\end{pmatrix}
=\begin{pmatrix}-1&3\\-1&2\end{pmatrix}\begin{pmatrix}13\\5\end{pmatrix}
=\begin{pmatrix}2\\-3\end{pmatrix}
\]
因此\(A=(2,-3)\).
\end{solution}
\end{problem}

\begin{problem}
已知直线 \(l\)：\(3x + 4y - 12 = 0\) 与圆 \(C\)：\(\begin{cases}
x = -1 + 2\cos \theta,\\
y = 2 + 2\sin \theta,
\end{cases}\) （\(\theta\) 为参数），试判断他们的公共点个数.

\begin{solution}
圆\(C\)的圆心为\(O_1=(-1,2)\)，半径为\(2\).圆心到直线\(l:3x+4y-12=0\)的距离为
\[
d=\frac{|3(-1)+4\times2-12|}{\sqrt{3^2+4^2}}=\frac{7}{5}<2
\]
因此直线与圆相交于两个不同的公共点，公共点个数为\(2\).
\end{solution}
\end{problem}

\begin{problem}
解不等式 \(|2x - 1| < |x| + 1\).

\begin{solution}
两边均为非负数，可以平方而不改变不等号方向，得
\[
(2x-1)^2<(|x|+1)^2
\quad\Longleftrightarrow\quad
3x^2-4x<2|x|
\]
当\(x\geqslant0\)时，\(|x|=x\)，所以
\(3x^2-6x<0\)，即\(0<x<2\).
当\(x<0\)时，\(|x|=-x\)，所以
\(3x^2-2x<0\)，即\(x(3x-2)<0\)，但此时两个因子均为负，乘积为正，无解.
综上，不等式的解集为\((0,2)\).
\end{solution}
\end{problem}
